% 4. Metodologija evaluacije
\section{Metodologija evaluacije}
\label{sec:metodologija}

EduGrader je projektovan kao platforma za ocenjivanje uz podršku veštačke inteligencije koja velike jezičke modele uključuje u tok provere znanja pod nadzorom nastavnika.
Primarni cilj sistema je unapređenje efikasnosti i doslednosti ocenjivanja uz očuvanje pedagoške transparentnosti i ljudske odgovornosti.
U ovom poglavlju opisana je metodologija evaluacije pouzdanosti ocena koje EduGrader generiše, skupovi podataka na kojima je sistem testiran i konfiguracije modela obuhvaćene eksperimentima.

\subsection{Metrike i evaluacija}
\label{sec:metrike}

Da bi se procenilo slaganje ocena generisanih veštačkom inteligencijom sa ljudskim sudom, izlazi sistema EduGrader upoređeni su sa ocenama koje su dodelila dva univerzitetska nastavnika.
Svaki nastavnik je ocenio drugačiji podskup od ukupno 686 odgovora (192 odnosno 485 odgovora).
Ovakva postavka omogućila je poređenje performansi sistema sa više ljudskih ocenjivača, umesto sa jednim jedinstvenim standardom.
Svaki odgovor ocenjen je prema standardnim kriterijumima ocenjivanja koji se koriste na datom kursu.
Ocene su dodeljivane na skali 0--10, na osnovu stepena konceptualne tačnosti i potpunosti odgovora.
Nastavnici su uzimali u obzir da li je student ispravno identifikovao ključne koncepte, da li nedostaju važni elementi i da li je rezonovanje delimično tačno, ali nepotpuno.
Odgovori koji demonstriraju ispravne osnovne ideje dobijali su delimične poene čak i kada nedostaju terminologija, struktura ili pojedini detalji, dok su suštinski netačni odgovori dobijali niske ocene ili nula poena.
Za svaki odgovor, ocena koju je dodelio nastavnik zadužen za kurs tretirana je kao referentna ocena i upoređena sa ocenom sistema pomoću metrika opisanih u nastavku.

Korišćena su četiri statistička pokazatelja koji obuhvataju komplementarne aspekte pouzdanosti ocenjivanja: Pirsonov koeficijent korelacije (PCC), kvadratno ponderisana Koenova kapa, odnos prosečnih ocena (engl. \emph{Mean Ratio}) i odnos korena srednje kvadratne greške i standardne devijacije (RMSE/STDDEV).

PCC~\cite{sedgwick2012pearson} meri jačinu linearne povezanosti između ocena koje su dodelili ljudi i ocena koje je dodelio sistem.
Koeficijent uzima vrednosti od $-1$ do $+1$, gde $+1$ označava savršenu pozitivnu linearnu vezu, $0$ odsustvo linearne veze, a $-1$ savršenu negativnu linearnu vezu.
Pre sprovođenja analize, raspodela razlika u ocenama ispitana je Šapiro--Vilkovim testom i vizuelnim pregledom Q--Q dijagrama.
Rezultati nisu ukazali na značajna odstupanja od normalnosti ($p > 0{,}05$), što opravdava upotrebu parametarske mere korelacije.
Važno je istaći da PCC meri linearnu povezanost, a ne apsolutno slaganje.
Iz tog razloga, PCC je dopunjen ponderisanom kapom kao merom slaganja i metrikama zasnovanim na grešci koje obuhvataju apsolutna odstupanja između ocena sistema i nastavnika.

Kao dopunu PCC-u, izveštavamo i kvadratno ponderisanu Koenovu kapu, koja je prikladna za ordinalne ocene.
Za razliku od PCC-a, koji meri linearnu povezanost, ponderisana kapa meri slaganje uzimajući u obzir uređenu prirodu skale ocenjivanja 0--10 i jače kažnjava veća neslaganja od manjih.

Odnos prosečnih ocena sažima opšte tendencije u ocenjivanju.
Razlike između prosečne ocene sistema i prosečne ocene nastavnika ukazuju na sistematsku blagost ili strogost, ali ne odražavaju usklađenost u variranju ocena.

Koren srednje kvadratne greške (RMSE)~\cite{hodson2022root} obuhvata prosečnu veličinu razlika između ljudskih ocena i ocena sistema; manja vrednost RMSE ukazuje na jače slaganje.
Odnos RMSE/STDDEV~\cite{chai2014root} meri koliko se ocene sistema slažu sa ljudskim ocenama u odnosu na prirodnu varijabilnost samih ljudskih ocena --- odnosno, da li su odstupanja sistema od ljudskog suda mala u poređenju sa tim koliko ljudske ocene uopšte variraju.
Vrednost odnosa RMSE/STDDEV blizu $0$ ukazuje na veoma visoko slaganje, vrednost oko $1$ znači da su odstupanja sistema uporediva sa tipičnom varijabilnošću ljudskog ocenjivanja, dok vrednost veća od $1$ sugeriše da se ocene sistema razlikuju od ljudskih više nego što se ljudi razlikuju međusobno.

\subsection{Prikupljanje podataka}
\label{sec:podaci}

Empirijska evaluacija sistema EduGrader sprovedena je na autentičnim odgovorima sa teorijskih ispita, prikupljenim tokom više ispitnih rokova na dva univerziteta u Srbiji.

Dobijeni skup podataka obuhvata \textbf{šest univerzitetskih kurseva} i pokriva više akademskih nivoa, uključujući osnovne i master studije.
Sadrži mešavinu konceptualnih pitanja i zadataka vezanih za programiranje, čime pruža raznovrstan i heterogen reper za evaluaciju sistema za automatsko ocenjivanje.
Kursevi obuhvaćeni skupom podataka su:

\begin{itemize}
    \item Programiranje 2 (Prog 2): kurs prve godine osnovnih studija, Matematički fakultet, Univerzitet u Beogradu;
    \item Tehničko i naučno pisanje (TSW): kurs prve godine osnovnih studija, Matematički fakultet, Univerzitet u Beogradu;
    \item Veb programiranje (WP): kurs četvrte godine osnovnih studija, Matematički fakultet, Univerzitet u Beogradu;
    \item Digitalno čuvanje podataka (DDS): kurs master studija, Univerzitet umetnosti;
    \item HTML5 i JavaScript za video-igre (HTML5): kurs master studija, Univerzitet umetnosti;
    \item Istraživanje podataka (DM): kurs treće godine osnovnih studija, Matematički fakultet, Univerzitet u Beogradu.
\end{itemize}

Uz svaki kurs, skup podataka sadrži studentske odgovore, ocene koje su dodelili nastavnici, referentne odgovore (kada su dostupni), napomene nastavnika i relevantne nastavne materijale korišćene u pripremi ispita.
Ukupno, skup sadrži \textbf{686 studentskih odgovora}, što predstavlja realističan reper za evaluaciju doslednosti ocenjivanja kroz različite discipline, akademske nivoe i formate pitanja.
U zavisnosti od kursa, odgovori uključuju \textbf{kratka tekstualna objašnjenja}, \textbf{duga tekstualna objašnjenja} i \textbf{odgovore u obliku koda}.
Tabela~\ref{tab:dataset} sumira strukturu skupa podataka, uključujući broj odgovora i njihove tipove za svaki kurs.

Radi bolje karakterizacije podataka, dodatno je prikazana raspodela tipova odgovora i prosečna dužina odgovora.
Ove statistike daju uvid u varijabilnost složenosti odgovora i pomažu u kontekstualizaciji težine automatskog ocenjivanja kroz različite formate pitanja.
Pored raspodele tipova odgovora, analizirana je i tipična dužina odgovora --- kako referentnih rešenja koja su pripremili nastavnici, tako i automatski generisanih referentnih odgovora i samih studentskih odgovora.
Ova informacija je relevantna jer dužina odgovora može uticati na težinu automatskog ocenjivanja, naročito kod dugih objašnjenja.

Tabela~\ref{tab:response_length} takođe pokazuje da su generisani referentni odgovori često znatno duži od odgovora koje su pripremili nastavnici, naročito kod pitanja gde se očekuju kratki odgovori.
Na više kurseva nastavnici očekuju sažete odgovore (npr.\ jedan izraz ili kratku konstataciju), dok odgovori koje generišu jezički modeli po pravilu uključuju i dodatna objašnjenja zašto je odgovor tačan.
Kao rezultat, generisani odgovori sadrže više kontekstualnih informacija nego što izvorno ispitno pitanje strogo zahteva.

\begin{table}[ht]
\centering
\caption{Struktura skupa podataka po kursevima (broj studentskih odgovora po tipu pitanja)}
\label{tab:dataset}
\resizebox{\columnwidth}{!}{%
\begin{tabular}{lcccc}
\toprule
\textbf{Kurs} &
\textbf{Odgovori} &
\textbf{Kratak tekst} &
\textbf{Dug tekst} &
\textbf{Kod} \\
\midrule
Programiranje 2 (Prog 2) & 480 & 380 & 0 & 100 \\
Tehničko i naučno pisanje (TSW) & 40 & 8 & 0 & 32 \\
Veb programiranje (WP) & 50 & 26 & 0 & 24 \\
Digitalno čuvanje podataka (DDS) & 10 & 10 & 0 & 0 \\
HTML5 i JavaScript za video-igre (HTML5) & 10 & 4 & 0 & 6 \\
Istraživanje podataka (DM) & 96 & 0 & 96 & 0 \\
\midrule
\textbf{Ukupno} & \textbf{686} & \textbf{428 (62{,}4\%)} & \textbf{96 (14\%)} & \textbf{162 (23{,}6\%)} \\
\bottomrule
\end{tabular}
}
\end{table}

\begin{table}[ht]
\centering
\caption{Prosečna dužina odgovora po tipu pitanja i izvoru odgovora}
\label{tab:response_length}
\resizebox{\columnwidth}{!}{%
\begin{tabular}{lcccccc}
\toprule
\textbf{Tip pitanja} &
\textbf{\shortstack{Odgovor nastavnika\\pros.\ reči}} &
\textbf{\shortstack{Odgovor nastavnika\\pros.\ znakova}} &
\textbf{\shortstack{Generisani odgovor\\pros.\ reči}} &
\textbf{\shortstack{Generisani odgovor\\pros.\ znakova}} &
\textbf{\shortstack{Studentski odgovor\\pros.\ reči}} &
\textbf{\shortstack{Studentski odgovor\\pros.\ znakova}} \\
\midrule
Kratak tekst & 20.9 & 130.5 & 201.7 & 1341.9 & 24.7 & 143.2 \\
Dug tekst    & 411.7 & 2867.7 & 479.9 & 3483.5 & 289.6 & 1889.3 \\
Kod          & 22.2 & 137.4 & 162.7 & 1144.3 & 15.9 & 102.5 \\
\bottomrule
\end{tabular}
}
\end{table}

\subsection{Modeli i konfiguracije}
\label{sec:modeli}

EduGrader koristi višemodelski pristup ocenjivanju koji uključuje i modele opšte namene i modele optimizovane za rezonovanje:
\begin{itemize}
    \item standardni modeli: GPT-4o, GPT-4o-mini, DeepSeek-Chat;
    \item modeli za rezonovanje: GPT-5.1, DeepSeek-Reasoner.
\end{itemize}
Svaki model evaluiran je u tri konfiguracije strogosti ocenjivanja: blagoj, neutralnoj i strogoj.
Ove konfiguracije realizovane su pažljivo projektovanim promptovima koji simuliraju različite filozofije ocenjivanja uobičajene u akademskoj praksi.
Na primer, stroga konfiguracija naglašava preciznost i potpunost (npr.\ „Ti si veoma strog profesor koji rigorozno ocenjuje studentske odgovore''), dok blaga konfiguracija nagrađuje delimičnu tačnost i konceptualno razumevanje.
Ovakav dizajn odražava nalaze prethodnih studija koje pokazuju da i ljudski ocenjivači značajno variraju u strogosti~\cite{jukiewicz2024future}.

Pored toga, EduGrader podržava dva komplementarna režima ocenjivanja:
\begin{itemize}
    \item PRA, u kome se studentski odgovori porede sa tačnim odgovorima i napomenama za ocenjivanje koje je pripremio nastavnik;
    \item GRA, u kome model najpre generiše referentni odgovor na osnovu priloženih nastavnih materijala, a zatim studentski odgovor ocenjuje u odnosu na taj generisani standard.
\end{itemize}
Ovakav dvorežimski dizajn omogućava evaluaciju kako u scenarijima sa dobro definisanim ključem za ocenjivanje, tako i u situacijama kada nastavnici ne obezbeđuju eksplicitna referentna rešenja.

Svi eksperimenti sprovedeni su na srpskom jeziku, uključujući ispitna pitanja, studentske odgovore i referentne odgovore.
Budući da je srpski jezik sa ograničenim resursima, slabo zastupljen u podacima na kojima se veliki jezički modeli obučavaju, ovo istraživanje proširuje oblast ocenjivanja uz podršku veštačke inteligencije izvan dominantnog engleskog konteksta.
